% Contenu partagé — corps des diapositives.
%
% Les mêmes environnements que sur papier, dans la structure propre au médium.
% C'est la vérification que l'API ne dépend pas du support.

%------------------------------------------------------------------
\begin{slide}{Les cinq familles à cadre}

\begin{theorem}{Cauchy--Lipschitz}{cauchyslide}
    Sous hypothèse de régularité locale, le problème de Cauchy admet une
    unique solution maximale.
\end{theorem}

\pause

\begin{definition}{Solution maximale}{}
    Une solution qui ne se prolonge pas.
\end{definition}

\pause

\begin{proposition}{}{}
    Une proposition.
\end{proposition}

\end{slide}

%------------------------------------------------------------------
\begin{slide}{Corollaire, conjecture, lemme}

\begin{corollary}{}{}
    Un corollaire.
\end{corollary}

\begin{conjecture}{}{}
    Une conjecture.
\end{conjecture}

\begin{lemma}
    Un lemme, composé au fil du texte.
\end{lemma}

\begin{example}
    Un exemple.
\end{example}

\begin{remark}
    Une remarque, avec son filet latéral.
\end{remark}

\end{slide}

%------------------------------------------------------------------
% La couleur d'en-tête change pour marquer un nouveau thème.
\begin{slide}[\ocotscolor{slide1}]{Une preuve sur plusieurs diapositives}

\begin{proposition}{}{}
    Pour tout $x \in \R$, $x^2 \ge 0$.
\end{proposition}

\begin{proofbegin}
    Premier morceau de la démonstration. Le carré final n'est pas encore posé.
\end{proofbegin}

\end{slide}

\begin{slide}[\ocotscolor{slide1}]{Suite de la preuve}

\begin{proofmiddle}
    Morceau intermédiaire.
\end{proofmiddle}

\pause

\begin{proofend}
    Dernier morceau. Le carré de fin arrive ici.
\end{proofend}

\end{slide}

%------------------------------------------------------------------
\begin{slide}[\ocotscolor{slide2}]{Blocs étiquetés}

\begin{assumption}
    Une hypothèse. L'étiquette ne change pas d'un \texttt{\textbackslash pause}
    à l'autre.
\end{assumption}

\pause

\begin{assumption}
    Une deuxième hypothèse.
\end{assumption}

\pause

\begin{openquestion}
    Un questionnement.
\end{openquestion}

\begin{difficulty}
    Une difficulté attendue.
\end{difficulty}

\end{slide}

%------------------------------------------------------------------
\begin{slide}[\ocotscolor{slide3}]{Mathématiques et mise en valeur}

Un \keyword{terme important}, puis \emphA{quatre} \emphB{couleurs}
\emphC{de} \emphD{thème}.

\medskip
Ensembles $\N$, $\Z$, $\R$, $\Rpos$, intervalles $\intervalleoo{a}{b}$,
normes $\norm{x}$, opérateurs $\rank A$ et $\trace A$.

\medskip
\[
    \tcbhighmath{\dot{x}(t) = A(t)\, x(t) + b(t)}
\]

\medskip
\begin{center}
\begin{tikzpicture}[xmin=-2, xmax=2, ymin=-1.2, ymax=1.2, scale=0.9]
    \grille
    \axes{$t$}{$x$}
    \draw[thick, color=\ocotscolor{emph-b}, domain=-2:2, samples=100]
        plot (\x, {sin(\x r)});
\end{tikzpicture}
\end{center}

\end{slide}

%------------------------------------------------------------------
\begin{slide}[\ocotscolor{slide4}]{Couleurs et liens}

Une \alert{mise en évidence projetée}, un lien externe avec
\myurl{https://www.example.com}, et un lien nommé vers la
slide~\hyperref[slide:exercise]{« exercice »}.

\begin{theorem}{Lien interne}{slide-link}
    Les liens internes doivent suivre la couleur du thème, comme les URL.
\end{theorem}

\end{slide}

%------------------------------------------------------------------
\begin{slide}[\ocotscolor{slide4}]{Un exercice}
\label{slide:exercise}

L'énoncé ci-dessous est copié tel quel depuis le polycopié.

\begin{exercise}[nosolution]
    Soit $E$ un espace vectoriel normé.
    \begin{question}
        Montrer que la boule $\ball(0,1)$ est convexe.
    \end{question}
    \begin{question}
        En déduire que $\norm{\cdot}$ est convexe.
    \end{question}
\end{exercise}

\end{slide}
